\section{Cơ sở lý thuyết}\label{sec:background}

\subsection{Kho hồ dữ liệu và Apache Iceberg}
Kho dữ liệu (\emph{data warehouse}) thường đưa dữ liệu vào các bảng được quản trị chặt để phục vụ phân tích ổn định. Ngược lại, hồ dữ liệu (\emph{data lake}) lưu nhiều loại tệp trên hạ tầng chi phí thấp nhưng có thể thiếu lớp quản trị bảng thống nhất. Kho hồ dữ liệu (\emph{Data Lakehouse}) kết hợp hai hướng: dữ liệu vẫn nằm ở định dạng mở, đồng thời được bổ sung siêu dữ liệu (\emph{metadata}), lịch sử phiên bản và ngữ nghĩa giao dịch ACID (\emph{Atomicity, Consistency, Isolation, Durability})~\parencite{armbrust2021lakehouse}.

Lưu trữ đối tượng (\emph{object storage}) truy cập dữ liệu theo khóa qua API HTTP thay vì đường dẫn tệp truyền thống. Trong Continux, MinIO cung cấp giao diện tương thích Amazon S3. Apache Iceberg là định dạng bảng (\emph{table format}), không phải bộ máy truy vấn: đặc tả này mô tả tệp nào thuộc bảng, ảnh chụp trạng thái (\emph{snapshot}) nào đang hiện hành và tệp kê khai (\emph{manifest}) nào theo dõi thay đổi~\parencite{icebergdocs}. Đích ghi (\emph{sink}) Iceberg có thể tạo tệp dữ liệu Parquet, tệp xóa theo giá trị khóa (\emph{equality-delete}), tệp xóa theo vị trí dòng (\emph{position-delete}) và siêu dữ liệu. Vì vậy, việc nhìn thấy các đối tượng trên MinIO chứng minh đường ghi đã hoạt động, không có nghĩa MinIO tự xử lý SQL.

\subsection{Xử lý luồng, trạng thái và khung nhìn vật lý}
Một sự kiện (\emph{event}) trong đề tài là bản ghi chuyến taxi được phát lại dưới dạng JSON. Xử lý luồng cập nhật kết quả khi sự kiện đến, thay vì chờ toàn bộ tệp hoàn tất như xử lý lô. Khung nhìn vật lý (\emph{materialized view}, MV) là kết quả truy vấn được hệ thống duy trì sẵn và cập nhật theo đầu vào. Với dữ liệu taxi, MV nhóm theo quận và khu vực, sau đó tính số chuyến, tổng cước và quãng đường trung bình.

Phép join với bảng tra cứu và phép tổng hợp theo vùng cần trạng thái (\emph{state}), tức thông tin trung gian được giữ qua nhiều sự kiện. Điểm kiểm tra (\emph{checkpoint}) là mốc lưu trạng thái có chủ đích để engine có cơ sở tiếp tục hoặc khôi phục xử lý sau gián đoạn. Continux dùng bucket \texttt{rw-checkpoint} trên MinIO cho trạng thái RisingWave.

\subsection{Broker tương thích Kafka và RisingWave}
Kafka định hình mô hình nhật ký phân tán cho xử lý luồng~\parencite{kreps2011kafka}. Chủ đề (\emph{topic}) là luồng thông điệp có tên; phân vùng (\emph{partition}) là phần chia phục vụ thứ tự cục bộ và xử lý song song; độ lệch (\emph{offset}) là vị trí của thông điệp trong phân vùng. Phía phát (\emph{producer}) ghi thông điệp, còn phía nhận (\emph{consumer}) đọc và ghi nhớ độ lệch đã xử lý.

Continux dùng Redpanda, broker tương thích Kafka API, để tách Vector khỏi RisingWave~\parencite{redpandadocs}. Topic \texttt{nyc-taxi-events} có ba phân vùng và một bản sao. Một bản sao nghĩa là đề tài không đánh giá khả năng chịu mất broker. RisingWave đọc topic, join với bảng tra cứu trên MinIO, duy trì MV và ghi kết quả ra Iceberg.

\subsection{Chuyển giao Xanh lam/Xanh lục}
Triển khai Xanh lam/Xanh lục (\emph{Blue/Green deployment}) duy trì hai phiên bản song song: Xanh lam là phiên bản đang phục vụ, Xanh lục là ứng viên thay thế. Chuyển giao (\emph{cutover}) đổi điểm truy cập công khai từ phiên bản cũ sang phiên bản mới. Trong Continux, đơn vị chuyển giao không phải pod Kubernetes mà là định nghĩa MV: \texttt{mv\_zone\_stats} là tên công khai ổn định, \texttt{mv\_zone\_stats\_blue} giữ cách tính nền và \texttt{mv\_zone\_stats\_green} áp dụng bộ lọc mới.

RisingWave hỗ trợ \texttt{ALTER MATERIALIZED VIEW ... SWAP WITH ...} để hoán đổi tên hai MV~\parencite{risingwaveSwap}. Tính nguyên tử (\emph{atomic}) ở đây chỉ áp dụng cho thao tác đổi tên nhìn từ truy vấn trong bộ máy: phía truy vấn không thấy trạng thái tên trung gian. Điều này không tự động chứng minh mọi quan hệ phụ thuộc phía sau cũng đổi theo. Tài liệu RisingWave yêu cầu xem xét riêng các thành phần phía sau (\emph{downstream}) như đích ghi khi thay đổi công việc xử lý luồng~\parencite{risingwaveAlterJobs}. Vì vậy kết luận của đề tài chỉ áp dụng cho tên truy vấn công khai.

\subsection{GitOps và khả năng quan sát}
OpenGitOps xác định bốn nguyên tắc cho trạng thái mong muốn: khai báo, lưu phiên bản trong Git, tự động kéo về và liên tục điều hòa sai lệch (\emph{reconcile})~\parencite{opengitops}. Continux dùng Argo CD theo mẫu ứng dụng-gốc-quản-lý-các-ứng-dụng-con (\emph{App-of-Apps})~\parencite{argocddocs}. Git giữ cấu hình dài hạn; việc scale Vector trong một lượt phát lại là thao tác đo có chủ đích.

Khả năng quan sát (\emph{observability}) là khả năng suy luận trạng thái nội tại từ tín hiệu bên ngoài như chỉ số, nhật ký và trạng thái khối lượng công việc. VictoriaMetrics lưu chuỗi thời gian~\parencite{victoriametricsdocs}; Grafana trực quan hóa chỉ số~\parencite{grafanadocs}; bộ xuất số đo (\emph{exporter}) của Continux chuyển kết quả SQL và kết quả chuyển giao thành số đo. Bảng điều khiển (\emph{dashboard}) hỗ trợ đọc hiện tượng, nhưng kết luận trong báo cáo luôn được đối chiếu với SQL và danh sách đối tượng MinIO.

\subsection{Ursa và nghiên cứu liên quan}
Ursa là bộ máy xử lý luồng tương thích Kafka và gắn trực tiếp với kho hồ dữ liệu, tập trung giảm chi phí sao chép và chi phí đầu nối (\emph{connector})~\parencite{merli2025ursa}. Continux không triển khai Ursa và không so sánh hiệu năng trực tiếp với Ursa. Đề tài khảo sát lớp vận hành phía trên: GitOps, quan sát hệ thống và thay đổi cách tính truy vấn công khai trên cụm nhỏ.

Flink hỗ trợ xử lý lô và xử lý luồng trong cùng một bộ máy~\parencite{carbone2015flink}; Kafka Streams và ksqlDB đại diện cho hướng xử lý gắn với hệ sinh thái Kafka. RisingWave được chọn vì giao diện SQL, MV và cú pháp hoán đổi phù hợp với hiện vật. Đây là định vị khái niệm, không phải phép đo hiệu năng đối đầu.
